% ESTRUCTURA DEL PROYECTO
%   main.tex                 Preámbulo (paquetes, portada, metadatos) y ensamblado
%   referencias.bib          Fuentes en formato BibTeX (biblatex-apa)
%   imagenes/                Logo institucional y demás figuras (Logo_UD.png)
%   secciones/               Contenido, un archivo por sección:
%     introduccion.tex
%     seccion1.tex
%     seccion2.tex
%     seccion3.tex
%     seccion4.tex
%     seccion5.tex
%     seccion6.tex
%     conclusiones.tex
%
% -----------------------------------------------------------------------
\documentclass[stu, 12pt, letterpaper, donotrepeattitle, floatsintext]{apa7}

\usepackage[T1]{fontenc}
\usepackage[utf8]{inputenc}
\usepackage[spanish]{babel}
\usepackage{csquotes}
\usepackage{newtxtext,newtxmath}
\usepackage{graphicx}
\graphicspath{{imagenes/}}
\usepackage{array}
\usepackage{booktabs}
\usepackage{longtable}
\usepackage{makecell}
\usepackage{tabularx}
\usepackage{setspace}
\usepackage[style=apa, backend=biber]{biblatex}
\DeclareLanguageMapping{spanish}{spanish-apa}
\addbibresource{referencias.bib}
\hypersetup{hidelinks}
\usepackage{pdflscape}
\usepackage{xcolor}
\usepackage{colortbl}
\usepackage{pgfgantt}


\newcolumntype{Y}{>{\raggedright\arraybackslash}X}

% Numeración de secciones y subsecciones: apa7 las deja sin numerar por defecto
% (secnumdepth = 0). Con 2 se obtiene "3" y "3.1", "3.2", ...; los
% subsubtemas y los \paragraph siguen sin número.
\setcounter{secnumdepth}{2}

% Ajuste menor: "s. f." con espacio (biblatex-apa lo deja como "s.f.").
\DefineBibliographyStrings{spanish}{%
	nodate = {s.~f\adddot},
}

\AtBeginDocument{%
	\renewcommand{\keywordname}{Palabras clave}%
	\renewcommand{\notelabel}{Nota}%
	\renewcommand{\tablename}{Tabla}% babel-spanish lo deja como "Cuadro"
}

\shorttitle{Fase 2 del Proyecto Final}
\title{%
	\includegraphics[width=1.9in]{Logo_UD.png}\\[0.6cm]
	Actividad de Aprendizaje No.~2.1: Fase 2 del Proyecto Final%
}
\authorsnames{Nicolás Guevara Herrán (20211180026), Jhojan Stiven Aragón Ramírez (20221020060), Daniel Orlando Villescas Mora (20221020065)}
\authorsaffiliations{{Maestría en Gestión y Seguridad de la Información, Facultad de Tecnología\\Universidad Distrital Francisco José de Caldas}}
\course{Arquitectura de Seguridad}
\professor{Camilo Enrique Rocha Calderon}
\duedate{22 de septiembre de 2026\\Bogotá, D.C.} % TODO: completar la fecha de entrega

\begin{document}
\maketitle
% Índice
\tableofcontents
\newpage

% Introducción
% Sin número: la numeración de secciones empieza en la sección 1 (Visión).
\section*{Introducción}
\addcontentsline{toc}{section}{Introducción}


% 1. Visión, Principios y Alcance de la Arquitectura
\section{Visión, Principios y Alcance de la Arquitectura}
% 2. Contexto de Negocio, Objetivos y Cumplimiento
\section{Contexto de Negocio, Objetivos y Cumplimiento}

% 3. Evaluación del Estado Actual y Análisis de Brechas
\section{Evaluación del Estado Actual y Análisis de Brechas}

% 4. Modelo de Amenazas, Activos y Evaluación de Riesgos
\section{Modelo de Amenazas, Activos y Evaluación de Riesgos}

Un modelo de amenazas describe qué protege la arquitectura, frente a qué amenazas y por qué medios. Parte del inventario de activos, caracteriza la superficie de ataque y valora el riesgo para decidir qué controles se necesitan y en qué orden. El modelo se apoya en la trazabilidad: cada control debe poder explicarse por un riesgo y cada riesgo por un activo del proceso.

En el Centro de Control Integrado (CCI), el modelo de amenazas y riesgos parte de los activos
críticos identificados durante la primera fase del proyecto \parencite{fase1cci2026} y mantiene la trazabilidad con los
riesgos, las políticas de seguridad, los mecanismos de protección y las capacidades arquitectónicas
desarrolladas posteriormente mediante SABSA, O-ESA y ArchiMate, y aplicadas en los ensayos y el modelo del caso \parencite{sabsaCCI2026,oesaCCI2026,modeloArchimateCCI2026}.

El análisis se concentra en los tres activos que soportan directamente el proceso de gestión y
despacho de flota: la telemetría de flota, la bitácora de decisiones del CCI y el Sistema de
Despacho. Sobre estos activos se mantienen los cuatro riesgos priorizados desde la primera fase \parencite{fase1cci2026}:
manipulación o suplantación de la telemetría (R1), alteración o eliminación de la bitácora
operativa (R2), acceso no autorizado al Sistema de Despacho (R3) e indisponibilidad del
Sistema de Despacho (R4).

Esta sección consolida la criticidad de los activos, caracteriza la superficie de ataque
mediante la taxonomía \textit{Data}, \textit{Method} y \textit{Channel} utilizada en O-ESA \parencite{oesa2011}, recupera la
valoración de riesgo definida previamente y relaciona cada riesgo con los dominios de seguridad
y las capacidades técnicas requeridas para alcanzar el estado objetivo de la arquitectura.


\subsection{Inventario y Criticidad de Activos}

El inventario de activos es el registro de la información, el software y la infraestructura que sostienen un proceso, junto con el valor que tienen para la organización. La criticidad se asigna por el efecto que tendría perder confidencialidad, integridad o disponibilidad en cada activo.

En el CCI, la identificación realizada en la primera fase \parencite{fase1cci2026} determinó tres elementos de criticidad
alta para la operación. La valoración de cada uno considera el impacto que tendría una afectación de
confidencialidad, integridad o disponibilidad sobre la regulación de la flota.

La Tabla~\ref{tab:activos_criticos_blueprint} resume esta clasificación.

\begin{table}[H]
	\centering
	\caption{Inventario y criticidad de los activos priorizados del CCI}
	\label{tab:activos_criticos_blueprint}
	
	\small
	\setlength{\tabcolsep}{4pt}
	\renewcommand{\arraystretch}{1.15}
	
	\begin{tabularx}{\textwidth}{
			>{\raggedright\arraybackslash}p{1.5cm}
			>{\raggedright\arraybackslash}X
			>{\centering\arraybackslash}p{2.0cm}
			>{\centering\arraybackslash}p{1.9cm}
			>{\centering\arraybackslash}p{2.0cm}
			>{\centering\arraybackslash}p{1.9cm}}
		
		\toprule
		
		\textbf{ID} &
		\textbf{Activo} &
		\shortstack{\textbf{Confiden-}\\\textbf{cialidad}} &
		\textbf{Integridad} &
		\shortstack{\textbf{Disponibi-}\\\textbf{lidad}} &
		\textbf{Criticidad} \\
		
		\midrule
		
		CCI-001 &
		Telemetría de flota &
		Media &
		Alta &
		Alta &
		Alta \\
		
		CCI-002 &
		Bitácora de decisiones del CCI &
		Media &
		Alta &
		Media &
		Alta \\
		
		CCI-003 &
		Sistema de Despacho &
		Media &
		Alta &
		Alta &
		Alta \\
		
		\bottomrule
		
	\end{tabularx}
	
	\vspace{2pt}
	
	\begin{minipage}{0.96\textwidth}
		\footnotesize
		\textit{Nota.} Elaboración propia con base en la identificación y valoración de activos
		realizada en la primera fase del proyecto \parencite{fase1cci2026}.
	\end{minipage}
\end{table}

La telemetría de flota es la fuente de información utilizada para conocer la
posición, velocidad y estado operacional de los vehículos. Su integridad y disponibilidad son
críticas porque los datos GPS/AVL alimentan el análisis operacional y
las decisiones posteriores de regulación.

La bitácora de decisiones del CCI conserva la evidencia de las acciones realizadas
durante la operación, incluyendo responsables, marcas de tiempo, instrucciones emitidas y
resultados. Su principal necesidad de protección se concentra en la integridad, ya que una
alteración o eliminación de registros afectaría la trazabilidad y la capacidad de auditoría.

El Sistema de Despacho es el principal activo de software mediante el cual se ejecutan las decisiones de regulación. Su integridad y su disponibilidad son críticas: una modificación no autorizada puede alterar las instrucciones operacionales, y una indisponibilidad puede impedir la aplicación de medidas ante retrasos, congestiones, bloqueos u otras novedades de la operación.


\subsection{Caracterización de la Superficie de Ataque}

La superficie de ataque es el conjunto de puntos por donde una amenaza puede entrar en contacto con un activo. Se descompone en los datos que pueden verse comprometidos, los métodos con que se materializa la amenaza y los canales físicos o lógicos por donde ocurre la interacción. Esa lectura evita tratar cada amenaza como un hecho aislado y permite ubicar el control en el punto de la arquitectura donde corresponde.

En el CCI, la caracterización sigue las tres dimensiones del análisis O-ESA \parencite{oesa2011,oesaCCI2026}: \textit{Data}, \textit{Method} y \textit{Channel}. \textit{Data} corresponde a la información o los elementos que pueden verse comprometidos; \textit{Method}, al mecanismo o técnica mediante la cual puede materializarse la amenaza; y \textit{Channel}, al canal físico o lógico a través del cual ocurre la interacción. Aplicada a los cuatro riesgos priorizados, esta descomposición muestra dónde se concentra la exposición de cada uno y qué tipo de control la reduce.

\begin{table}[H]
	\centering
	\caption{Caracterización de la superficie de ataque de los riesgos priorizados}
	\label{tab:superficie_ataque}
	\scriptsize
	\begin{tabularx}{\textwidth}{
			>{\raggedright\arraybackslash}p{1.15cm}
			>{\raggedright\arraybackslash}p{2.7cm}
			>{\raggedright\arraybackslash}X
			>{\raggedright\arraybackslash}X
			>{\raggedright\arraybackslash}X}
		\toprule
		\textbf{Riesgo} &
		\textbf{Activo principal} &
		\textbf{Data} &
		\textbf{Method} &
		\textbf{Channel} \\
		\midrule
		
		R1 &
		Telemetría de flota &
		Coordenadas GPS/AVL, velocidad, posición y estado operacional de los vehículos. &
		\textit{Spoofing}, manipulación o suplantación de mensajes de telemetría. &
		Comunicación 4G/5G vehículo--CCI y canal de ingestión hacia el bus de eventos. \\
		
		\addlinespace
		
		R2 &
		Bitácora de decisiones del CCI &
		Registros de decisiones, responsables, marcas de tiempo, instrucciones ejecutadas y resultados. &
		Modificación, borrado o truncamiento de registros y abuso de privilegios. &
		Sistema de Despacho, base de datos y mecanismos administrativos de almacenamiento. \\
		
		\addlinespace
		
		R3 &
		Sistema de Despacho &
		Identidades, credenciales, perfiles, roles y órdenes de despacho. &
		Robo o uso compartido de credenciales, abuso de privilegios y escalamiento de permisos. &
		Consolas del CCI, infraestructura SSO y aplicaciones que conforman el Sistema de Despacho. \\
		
		\addlinespace
		
		R4 &
		Sistema de Despacho e infraestructura tecnológica &
		Estado y disponibilidad del servicio, sesiones y flujo operacional. &
		Falla técnica, denegación de servicio, saturación o indisponibilidad de un componente crítico. &
		Infraestructura local, enlaces de telecomunicaciones, red del CCI y servicios desplegados en la nube. \\
		
		\bottomrule
	\end{tabularx}
	
	\begin{minipage}{0.96\textwidth}
		\footnotesize
		\textit{Nota.} Elaboración propia a partir de la matriz de trazabilidad de riesgos,
		superficie de ataque y servicios de seguridad desarrollada mediante O-ESA \parencite{oesa2011,oesaCCI2026}.
	\end{minipage}
\end{table}

En R1, la superficie de ataque se encuentra principalmente en la frontera entre un
entorno móvil no confiable y los sistemas internos del CCI. La información se origina en las
unidades ITS instaladas en los vehículos y posteriormente atraviesa las redes móviles y la
pasarela de ingestión antes de incorporarse a los sistemas que soportan la regulación de flota.
Por eso, la arquitectura objetivo requiere validar tanto la identidad del dispositivo como
la autenticidad, integridad y coherencia física de cada mensaje antes de incorporarlo a la fuente
autoritativa de datos.

Para R2, la superficie de ataque se concentra en los mecanismos mediante los cuales se generan, almacenan y administran los registros operacionales. El riesgo incluye ataques externos, abuso de privilegios y modificaciones posteriores sobre los registros. En consecuencia, la protección debe garantizar que la evidencia se capture automáticamente, que no dependa de la voluntad del operador y que cualquier alteración posterior pueda demostrarse.

En R3, los principales puntos de exposición son las identidades, credenciales, roles y permisos utilizados para interactuar con las consolas y con el Sistema de Despacho. El riesgo se materializa cuando una identidad no autorizada logra ejecutar acciones de regulación o cuando un usuario legítimo dispone de privilegios superiores a los necesarios para su función.

R4 presenta una superficie de ataque distribuida entre la infraestructura local, las comunicaciones y los servicios tecnológicos de los cuales depende el despacho. Además de ataques de denegación de servicio, deben contemplarse fallas de infraestructura, saturación y pérdida de componentes críticos, porque cualquiera de estos eventos puede afectar la disponibilidad de la operación.


\subsection{Evaluación de Riesgos Residuales}

El riesgo inherente es el nivel de exposición de un activo antes de aplicar controles; el riesgo residual es el que permanece después de tratarlo. La diferencia entre ambos mide el efecto real de los controles y permite decidir si el riesgo resultante cabe dentro del apetito de la organización o exige tratamiento adicional.

En el CCI, la valoración de la primera fase \parencite{fase1cci2026} estableció el nivel inherente de los
riesgos asociados a los activos críticos. El Anexo B consolidó los controles definidos para su tratamiento y calculó la valoración residual después de considerar su efecto sobre la probabilidad o el impacto.

Para mantener la trazabilidad con el resto del Blueprint, en esta sección se conserva la
numeración de riesgos utilizada a lo largo del proyecto: R1 para la manipulación de telemetría,
R2 para la alteración de la bitácora de decisiones, R3 para el acceso no autorizado al Sistema
de Despacho y R4 para la indisponibilidad del Sistema de Despacho.

La matriz del Anexo B utiliza una numeración diferente. En ella, la
indisponibilidad aparece como R1, la manipulación de telemetría como R2, el riesgo sobre la
integridad y no repudio de la bitácora como R3 y el acceso no autorizado al Sistema de
Despacho como R4. Por esta razón, los valores residuales se trasladan al Blueprint según la
descripción y el activo asociado a cada riesgo, y no únicamente según el identificador numérico.

La Tabla~\ref{tab:correspondencia_riesgos} presenta esta correspondencia.

\begin{table}[H]
	\centering
	\caption{Correspondencia entre los riesgos del Blueprint y el Anexo B}
	\label{tab:correspondencia_riesgos}
	\small
	\setlength{\tabcolsep}{5pt}
	\renewcommand{\arraystretch}{1.15}
	
	\begin{tabularx}{\textwidth}{
			>{\centering\arraybackslash}p{1.6cm}
			>{\raggedright\arraybackslash}X
			>{\centering\arraybackslash}p{2.0cm}}
		
		\toprule
		\textbf{ID Blueprint} &
		\textbf{Riesgo} &
		\textbf{ID Anexo B} \\
		\midrule
		
		R1 &
		Manipulación o suplantación de la telemetría de flota &
		R2 \\
		
		R2 &
		Alteración de la bitácora de decisiones del CCI &
		R3 \\
		
		R3 &
		Acceso no autorizado al Sistema de Despacho &
		R4 \\
		
		R4 &
		Indisponibilidad del Sistema de Despacho &
		R1 \\
		
		\bottomrule
	\end{tabularx}
	
	\begin{minipage}{0.96\textwidth}
		\footnotesize
		\textit{Nota.} Se conserva la identificación utilizada en el Blueprint y se establece su
		correspondencia con la matriz de riesgos del Anexo B.
	\end{minipage}
\end{table}


\subsubsection{Controles aplicados}

Los controles registrados en el Anexo B se orientan al tratamiento de los cuatro riesgos
priorizados y tienen como estrategia común la reducción del riesgo.

Para R1, asociado a la manipulación o suplantación de telemetría, se plantean controles de validación criptográfica, integridad y monitoreo continuo sobre los componentes ITS para detectar alteraciones en las señales recibidas desde la flota. Para R2, enfocado en la integridad y trazabilidad de la bitácora de decisiones, se fortalecen los accesos privilegiados mediante MFA, segregación de funciones y registros de auditoría que permitan garantizar el no repudio.

Para R3, relacionado con el acceso no autorizado al Sistema de Despacho, se implementan perfiles RBAC diferenciados, MFA en las consolas del CCI, auditoría de las órdenes emitidas y mecanismos de alertamiento ante actividades no autorizadas. Finalmente, para R4, correspondiente a la indisponibilidad, se propone un plan de recuperación ante desastres acompañado de mecanismos de redundancia, un sitio alterno de contingencia y pruebas periódicas de recuperación.


\subsubsection{Valoración residual}

La aplicación de los controles modifica la exposición de los riesgos de manera diferente. En
los riesgos R1, R2 y R3, los controles preventivos reducen la probabilidad residual desde el
nivel Muy Alto hasta el nivel Medio, correspondiente al 50\,\%. El impacto residual permanece
en el nivel Mayor, con un valor del 80\,\%.

En R4, asociado a la disponibilidad, el control definido tiene carácter principalmente
correctivo. Por esta razón, la probabilidad residual permanece en Muy Alta (100\,\%), mientras
que el impacto se reduce al 65\,\%, conservándose dentro de la categoría Mayor.

La severidad residual presentada en la Tabla~\ref{tab:riesgo_residual_blueprint} corresponde
al producto entre la probabilidad residual y el impacto residual. De acuerdo con la matriz del
Anexo B, los cuatro riesgos permanecen en una zona de riesgo final Alta, por lo cual
la estrategia de tratamiento definida para todos ellos continúa siendo Reducir.

\begin{table}[H]
	\centering
	\caption{Consolidación del riesgo residual de los riesgos priorizados}
	\label{tab:riesgo_residual_blueprint}
	\small
	\setlength{\tabcolsep}{3.5pt}
	\renewcommand{\arraystretch}{1.2}
	
	\begin{tabularx}{\textwidth}{
			>{\centering\arraybackslash}p{0.8cm}
			>{\raggedright\arraybackslash}X
			>{\centering\arraybackslash}p{1.8cm}
			>{\centering\arraybackslash}p{1.8cm}
			>{\centering\arraybackslash}p{1.55cm}
			>{\centering\arraybackslash}p{1.4cm}
			>{\centering\arraybackslash}p{2.2cm}}
		
		\toprule
		\textbf{ID} &
		\textbf{Riesgo} &
		\shortstack{\textbf{Prob.}\\\textbf{residual}} &
		\shortstack{\textbf{Impacto}\\\textbf{residual}} &
		\shortstack{\textbf{Severidad}\\\textbf{residual}} &
		\shortstack{\textbf{Nivel}\\\textbf{residual}} &
		\textbf{Tratamiento} \\
		\midrule
		
		R1 &
		Manipulación o suplantación de telemetría &
		Media (50\,\%) &
		Mayor (80\,\%) &
		40\,\% &
		Alto &
		Reducir \\
		
		R2 &
		Alteración de la bitácora de decisiones &
		Media (50\,\%) &
		Mayor (80\,\%) &
		40\,\% &
		Alto &
		Reducir \\
		
		R3 &
		Acceso no autorizado al Sistema de Despacho &
		Media (50\,\%) &
		Mayor (80\,\%) &
		40\,\% &
		Alto &
		Reducir \\
		
		R4 &
		Indisponibilidad del Sistema de Despacho &
		Muy Alta (100\,\%) &
		Mayor (65\,\%) &
		65\,\% &
		Alto &
		Reducir \\
		
		\bottomrule
	\end{tabularx}
	
	\begin{minipage}{0.96\textwidth}
		\footnotesize
		\textit{Nota.} Elaboración propia con base en la valoración residual consignada en el
		Anexo B. Los identificadores se presentan según la nomenclatura utilizada en el Blueprint;
		la correspondencia con los identificadores del Anexo B se presenta en la
		Tabla~\ref{tab:correspondencia_riesgos}.
	\end{minipage}
\end{table}

Los resultados muestran que los controles definidos reducen de manera importante la
probabilidad de materialización de los riesgos asociados a telemetría, trazabilidad y control
de acceso. Sin embargo, ninguno de los cuatro riesgos desciende todavía a una zona inferior a
Alta.

R4 conserva una probabilidad Muy Alta porque la exposición del Sistema de Despacho y de la infraestructura del CCI es permanente. La estrategia de continuidad acepta que una falla puede ocurrir y busca limitar sus consecuencias mediante redundancia, contingencia y recuperación.

En consecuencia, los cuatro riesgos requieren mantener una estrategia de reducción y
seguimiento continuo. La permanencia del nivel residual Alto también justifica las capacidades
de seguridad planteadas en la arquitectura objetivo y su incorporación en los estados de
transición y en la hoja de ruta desarrollada posteriormente.


\subsection{Matriz de Dominios de Seguridad}

La matriz de dominios de seguridad agrupa los requerimientos de la arquitectura objetivo en áreas de capacidad. Cada dominio reúne los controles que comparten un mismo riesgo y una misma forma de verificarse, de modo que la implementación pueda planificarse y evaluarse por separado.

En el CCI, los riesgos priorizados, el análisis de brechas y las capacidades modeladas en ArchiMate \parencite{archimate2022,modeloArchimateCCI2026} dan lugar a cuatro dominios de intervención. Los dominios se definen por los requerimientos técnicos y de control que debe satisfacer la arquitectura objetivo. Las tecnologías se presentan después como mecanismos de implementación de esos requerimientos.

\begin{table}[H]
	\centering
	\caption{Matriz de dominios de seguridad para la arquitectura objetivo del CCI}
	\label{tab:matriz_dominios_seguridad}
	\scriptsize
	\begin{tabularx}{\textwidth}{
			>{\raggedright\arraybackslash}p{2.35cm}
			>{\raggedright\arraybackslash}p{1.15cm}
			>{\raggedright\arraybackslash}X
			>{\raggedright\arraybackslash}X}
		\toprule
		\textbf{Dominio de seguridad} &
		\textbf{Riesgo} &
		\textbf{Requerimientos técnicos y de control} &
		\textbf{Componentes / tecnologías de implementación} \\
		\midrule
		
		Gestión de identidad y acceso &
		R3 &
		Toda interacción con funciones críticas del Sistema de Despacho debe asociarse a una identidad
		individual. Se requiere autenticación multifactor obligatoria, autorización basada en roles,
		mínimo privilegio, segregación de funciones y evaluación centralizada de políticas antes de
		permitir acciones de regulación. &
		Active Directory y ADFS para identidad federada y SSO; MFA; perfiles RBAC; repositorio de
		políticas; motor PDP; PEP integrados en consolas, aplicaciones e interfaces; monitoreo de
		eventos de autenticación y autorización. \\
		
		\addlinespace
		
		Seguridad e integridad de telemetría &
		R1 &
		Cada unidad ITS debe disponer de una identidad verificable y cada trama debe demostrar
		autenticidad e integridad antes de incorporarse al flujo operacional. Además, deben aplicarse
		validaciones automáticas de plausibilidad y cuarentena para mensajes inválidos o anómalos. &
		Certificados X.509 y/o mecanismos HMAC; autenticación mutua; protección criptográfica de
		comunicaciones; pasarela de validación; reglas de plausibilidad; bus de eventos/RabbitMQ;
		Oracle Spatial como repositorio geoespacial autoritativo. \\
		
		\addlinespace
		
		Trazabilidad, auditoría e inmutabilidad &
		R2 &
		Todas las decisiones críticas deben registrarse automáticamente y conservar la relación entre
		evento de origen, responsable, instrucción, marca de tiempo y resultado. Los registros deben
		protegerse contra modificación o eliminación y deben existir verificaciones periódicas de su
		integridad. &
		Bitácora centralizada; almacenamiento WORM o de solo adición; sellado y encadenamiento de
		hashes; sincronización temporal mediante NTP; FortiAnalyzer/FortiSIEM para correlación y
		alertamiento; procedimientos de cadena de custodia digital. \\
		
		\addlinespace
		
		Disponibilidad, resiliencia y continuidad &
		R4 &
		La arquitectura debe eliminar puntos únicos de falla, separar dominios de confianza, disponer
		de redundancia para los componentes críticos y garantizar mecanismos de recuperación y
		conmutación que permitan restablecer la operación dentro de los objetivos de continuidad
		establecidos. &
		Infraestructura local del CCI y Google Cloud Platform; clústeres y firewalls Fortinet;
		segmentación de red; balanceadores de carga; replicación de servicios y datos; failover
		automatizado; monitoreo de disponibilidad; DRP y BCP. \\
		
		\bottomrule
	\end{tabularx}
	
	\begin{minipage}{0.96\textwidth}
		\footnotesize
		\textit{Nota.} Elaboración propia a partir del análisis de brechas, las capacidades modeladas
		en ArchiMate \parencite{archimate2022,modeloArchimateCCI2026} y la arquitectura tecnológica de seguridad definida mediante O-ESA \parencite{oesa2011,oesaCCI2026}.
	\end{minipage}
\end{table}

La matriz mantiene una relación directa entre los riesgos identificados y las capacidades
que deben incorporarse a la arquitectura objetivo. Los componentes tecnológicos implementan
requerimientos derivados del riesgo y de las necesidades operacionales del CCI.

Estos cuatro dominios son la base de la arquitectura de transición de la Sección~5. Las
capacidades de identidad y acceso, seguridad de telemetría, trazabilidad de decisiones y continuidad
operativa deben evolucionar desde el estado actual hacia los estados de transición que allí se describen.
% 5. Arquitectura de Transición y Hoja de Ruta
\section{Arquitectura de Transición y Hoja de Ruta}

% 6. Métricas
\section{Métricas}

La medición de la arquitectura de seguridad permite verificar que los controles y capacidades
definidos en las secciones anteriores produzcan los resultados esperados sobre la operación del CCI. Para ello se retoma el perfil de atributos de negocio
desarrollado mediante SABSA, en el cual cada requerimiento de seguridad se encuentra asociado
con indicadores de desempeño (\textit{Key Performance Indicators}, KPI) e indicadores de
riesgo (\textit{Key Risk Indicators}, KRI).

Los KPI representan las metas de desempeño que debe alcanzar la arquitectura y permiten
evaluar la efectividad de los controles implementados. Los KRI funcionan como mecanismos de
alerta temprana y establecen umbrales amarillo y rojo que facilitan identificar desviaciones antes
de que estas generen un nivel de exposición no aceptable.

Esta estructura permite relacionar directamente las métricas con los riesgos R1 a R4, así como
con los dominios de seguridad, servicios O-ESA y capacidades arquitectónicas implementadas
durante las diferentes fases de la hoja de ruta.


\subsection{Indicadores y Métricas de Desempeño (KPI / KRI)}

La Tabla~\ref{tab:kpi_kri_1} y la Tabla~\ref{tab:kpi_kri_2} consolidan los indicadores
definidos para evaluar la integridad y oportunidad de la telemetría, el control de acceso,
la trazabilidad de las decisiones, la disponibilidad y recuperación del servicio, el monitoreo
de seguridad, la protección de datos personales y el cumplimiento normativo.


\begin{table}[H]
	\centering
	\caption{Indicadores KPI y KRI de la arquitectura de seguridad del CCI (BA1--BA6)}
	\label{tab:kpi_kri_1}
	
	\begingroup
	\scriptsize
	\setstretch{1}
	\setlength{\tabcolsep}{3pt}
	\renewcommand{\arraystretch}{1.12}
	
	\begin{tabularx}{\textwidth}{
			>{\centering\arraybackslash}p{0.8cm}
			>{\raggedright\arraybackslash}p{2.0cm}
			>{\centering\arraybackslash}p{1.4cm}
			Y
			Y}
		
		\toprule
		\textbf{ID} &
		\textbf{Atributo} &
		\textbf{Riesgo} &
		\textbf{KPI (meta)} &
		\textbf{KRI (amarillo / rojo)} \\
		\midrule
		
		BA1 &
		Integridad asegurada &
		R1 &
		Al menos 99,5\% de los mensajes GPS/AVL con autenticidad e integridad verificadas. &
		Registros rechazados o anómalos por cada 100.000: amarillo $>$ 50; rojo $>$ 200. \\
		
		\addlinespace[2pt]
		
		BA2 &
		Exacta y actual &
		R1 &
		Latencia de posición p95 $\leq$ 30 s. &
		Buses sin reporte vigente por más de 2 min: amarillo $>$ 2\%; rojo $>$ 5\%. \\
		
		\addlinespace[2pt]
		
		BA3 &
		Autenticada y autorizada &
		R3 &
		100\% de cuentas de consolas críticas con MFA y 100\% de revisiones trimestrales
		de permisos ejecutadas. &
		Accesos fuera de perfil u horario por 1.000 sesiones: amarillo $>$ 1; rojo $>$ 5. \\
		
		\addlinespace[2pt]
		
		BA4 &
		Segregación de funciones &
		R2, R3 &
		Cero usuarios con roles incompatibles. &
		Excepciones de segregación abiertas: amarillo $\geq$ 1; rojo $\geq$ 3. \\
		
		\addlinespace[2pt]
		
		BA5 &
		Trazable y responsable &
		R2 &
		Al menos 99,9\% de decisiones con registro completo. &
		Decisiones sin registro completo por 1.000: amarillo $>$ 1; rojo $>$ 5. \\
		
		\addlinespace[2pt]
		
		BA6 &
		Auditable y no repudiable &
		R2 &
		100\% de verificaciones periódicas de integridad del repositorio WORM aprobadas. &
		Intentos de modificar o borrar registros: amarillo = 1 bloqueado; rojo = cualquier
		alteración efectiva. \\
		
		\bottomrule
	\end{tabularx}
	
	\endgroup
\end{table}


\begin{table}[H]
	\centering
	\caption{Indicadores KPI y KRI de la arquitectura de seguridad del CCI (BA7--BA11)}
	\label{tab:kpi_kri_2}
	
	\begingroup
	\scriptsize
	\setstretch{1}
	\setlength{\tabcolsep}{3pt}
	\renewcommand{\arraystretch}{1.12}
	
	\begin{tabularx}{\textwidth}{
			>{\centering\arraybackslash}p{0.8cm}
			>{\raggedright\arraybackslash}p{2.0cm}
			>{\centering\arraybackslash}p{1.4cm}
			Y
			Y}
		
		\toprule
		\textbf{ID} &
		\textbf{Atributo} &
		\textbf{Riesgo} &
		\textbf{KPI (meta)} &
		\textbf{KRI (amarillo / rojo)} \\
		\midrule
		
		BA7 &
		Disponible &
		R4 &
		Disponibilidad mensual $\geq$ 99,9\%. &
		Minutos de indisponibilidad en el mes: amarillo $>$ 20; rojo $>$ 43. \\
		
		\addlinespace[2pt]
		
		BA8 &
		Continua y recuperable &
		R4 &
		RTO $\leq$ 15 min y 100\% de pruebas semestrales de conmutación exitosas. &
		Días desde la última prueba exitosa: amarillo $>$ 150; rojo $>$ 180. \\
		
		\addlinespace[2pt]
		
		BA9 &
		Detectable y monitoreada &
		R1--R4 &
		Tiempo medio de detección $\leq$ 5 min en eventos críticos y al menos 95\% de
		activos críticos reportando al SIEM. &
		Alertas críticas sin atender por más de 15 min: amarillo $\geq$ 1; rojo $\geq$ 5. \\
		
		\addlinespace[2pt]
		
		BA10 &
		Confidencial y privada &
		Transversal &
		100\% de indicadores y tableros sin datos personales identificables innecesarios. &
		Accesos a datos personales sin justificación: amarillo $\geq$ 1; rojo $\geq$ 3. \\
		
		\addlinespace[2pt]
		
		BA11 &
		Cumplimiento &
		Transversal &
		Al menos 90\% de controles aplicables con evidencia vigente. &
		Hallazgos de auditoría vencidos: amarillo $\geq$ 1; rojo $\geq$ 3. \\
		
		\bottomrule
	\end{tabularx}
	
	\endgroup
\end{table}


Los indicadores permiten evaluar diferentes dimensiones de la arquitectura. BA1 y BA2
verifican la confiabilidad y oportunidad de la telemetría utilizada por el Sistema de Despacho;
BA3 y BA4 evalúan la efectividad del gobierno de identidades y la segregación de funciones;
BA5 y BA6 miden la trazabilidad, integridad y capacidad probatoria de la bitácora; mientras que
BA7 y BA8 permiten comprobar la disponibilidad y capacidad de recuperación del CCI.

De manera transversal, BA9 evalúa la capacidad de detección y monitoreo de eventos sobre los
activos críticos, mientras que BA10 y BA11 permiten verificar respectivamente la protección de
datos personales y el cumplimiento de los controles normativos aplicables.

El seguimiento de estos indicadores debe integrarse con los mecanismos de monitoreo y
auditoría definidos en la arquitectura objetivo. Los KRI que alcancen el nivel amarillo deben
generar una revisión preventiva del control asociado, mientras que un umbral rojo requiere
escalamiento y tratamiento prioritario por parte de los responsables de seguridad y operación.

Finalmente, las métricas deben utilizarse como insumo para revisar periódicamente el riesgo
residual de R1 a R4. De esta forma, la arquitectura no se considera efectiva únicamente por la
implementación de los mecanismos técnicos definidos en la hoja de ruta, sino por su capacidad
demostrable para mantener la integridad de los datos, limitar accesos no autorizados, preservar
la trazabilidad y sostener la disponibilidad de los servicios críticos del CCI.

% Conclusiones
\section{Conclusiones}

Un Blueprint de Arquitectura de Seguridad reúne los principales elementos que deben considerarse para
responder a las necesidades de seguridad de una organización, precisa el propósito de cada uno y
muestra cómo se articulan. El caso del Centro de Control Integrado (CCI) de TransMilenio~S.A. permite
recorrer esos elementos en una aplicación práctica, desde la visión, los principios y el alcance de la
arquitectura hasta el análisis de la situación actual, la identificación de brechas, el diseño de la
arquitectura objetivo y la planificación de su implementación.

El análisis del CCI mostró que la seguridad es crítica en un entorno que centraliza procesos de
supervisión y regulación de la flota. Se identificaron riesgos relacionados con la manipulación de la
telemetría, los accesos no autorizados, la trazabilidad de las decisiones operativas y la
disponibilidad de los servicios. Aunque existen controles como la autenticación, las herramientas de
seguridad perimetral y el almacenamiento inmutable WORM, su integración, cobertura y efectividad aún
pueden mejorar.

La arquitectura objetivo (TO-BE) se organiza en cuatro dominios principales: gobierno de identidades
y control de acceso, integridad de la telemetría, auditoría y trazabilidad de las decisiones, y
continuidad operativa. Los dominios organizan los controles según los riesgos identificados y
resumen las capacidades que necesita el CCI para operar de forma confiable, controlada y resiliente.

Las arquitecturas de transición y la hoja de ruta llevan esa visión a una implementación progresiva,
organizada en horizontes de tiempo. Esa secuencia permite priorizar las acciones: primero el
fortalecimiento de los mecanismos de acceso, después la validación e integridad de la información, y
al final las capacidades de alta disponibilidad y recuperación ante incidentes.

SABSA y O-ESA, junto con los modelos y las vistas de ArchiMate, permiten relacionar las necesidades
del negocio con los riesgos, los controles, los servicios y los componentes tecnológicos de la
arquitectura. Los indicadores y las métricas, por su parte, sirven para evaluar el cumplimiento de
los objetivos de seguridad y hacer seguimiento a la efectividad de las iniciativas propuestas.

El caso del CCI muestra que un Blueprint de Arquitectura de Seguridad es una guía que conecta las
necesidades de la organización con los riesgos, las capacidades de seguridad y las acciones
necesarias para alcanzar un estado objetivo. Los controles y las tecnologías se eligen dentro de esa
guía, después de definir los riesgos y las capacidades. Su elaboración exige considerar en conjunto
los aspectos estratégicos, organizacionales, tecnológicos y operativos, y traducir el diseño en
decisiones y procesos de implementación.

% La lista de referencias se genera automáticamente con las entradas
% que se citen en el texto (\parencite / \textcite) desde referencias.bib.
% Si aún no hay citas, esta sección no imprime nada (aviso "Empty bibliography").
\printbibliography

\end{document}
